\section{Introduction}
\label{sec:introduction}

Cumulative incidence is a common observation in epidemic modelling when infection times are hidden and only counts of new cases are recorded in successive windows. \citet{whitaker2025} take a time-discretised susceptible-exposed-infected-removed model, and its susceptible-infected-removed reduction, as the latent transmission process, and they assimilate noisy incidence with a particle filter. The same construction matters more widely for partially observed jump processes, where a conditioned proposal is used because forward simulation ignores the next observation \citep{golightly2015,golightly2019}.

The inferential obstacle is an intractable observed-data likelihood. The data are incomplete and discrete, the reporting model is noisy, and the latent incidence between two observation times is a path of counts. A bootstrap filter would draw that path from the unconditioned transition and weight by the observation alone. \citet[p.~8]{whitaker2025} give two reasons not to stop there. Resampling static parameters collapses them, and forward simulation of incidence ignores the next count.

\citet{whitaker2025} meet the second difficulty with a conditioned hazard. Particles are propagated by replacing the Poisson rate with a hazard derived from a Gaussian approximation of the cumulative incidence between the current time and the next observation, given that observation. The particle is then reweighted by the discretised transition and the observation mass function. The limitations section of that paper states that a further assumption in the derivation is linearity of the latent process over an interval no longer than the observation interval. When the epidemic is sparsely observed, so that cumulative incidence does not evolve linearly, ``the conditioned hazard is likely to result in trajectories that are inconsistent with the true conditioned process'' \citep[p.~25]{whitaker2025}.

That paragraph gives no numerical spacing at which the failure begins, and it does not say that the importance weight then targets a different distribution. It names recent work on sampling conditioned jump processes, in particular \citet{golightly2019} and \citet{corstanje2023}, as relevant and open to adaptation. Neither algorithm is printed there as a replacement inside the filter.

Our contribution is to replace the Poisson means inside that proposal and to measure the effect on the weights. We treat the constant-hazard step in the 2025 derivation as the linearity assumption named on that page, so that the inconsistency described there is a property of the proposal. We report the effective sample size of the incidence weights on the synthetic susceptible-infected-removed mechanism of \citet{whitaker2025}, at the published spacing and at two coarser spacings fixed in advance. Equation (19) of \citet{golightly2019} and Theorem 4.2 of \citet{corstanje2023} describe proposals for objects other than the incidence increments in the filter. Section~\ref{sec:residual} states the replacement we use: the increment of a restarted incidence linear-noise path, plus a conditional residual mean for infections. The density in the denominator remains a product of Poisson probabilities for those incidence increments.

The numerical illustration uses the synthetic mechanism in Section 4.1 of \citet{whitaker2025}. That paper also discusses weekly Ebola counts and weekly COVID-19 counts. The comparison in this paper is confined to the synthetic mechanism.

The remainder of this paper is organised as follows. Section~\ref{sec:related} compares related guides. Section~\ref{sec:filter} restates the 2025 filter and the linear incidence guide. Section~\ref{sec:diagnostic} states the diagnostic. Section~\ref{sec:residual} states the residual proposal. Section~\ref{sec:simulation} reports both comparisons, and Section~\ref{sec:discussion} discusses their scope.
